% =====================================================================
%  Palette colori
%  Ricavata dal PDF della tesi del 2022 (legacy/tesi-magistrale/), leggendo
%  i valori nei content stream e nelle annotazioni dei link:
%    - rosso Sapienza RGB(130,36,51): frontespizio (definito dalla classe
%      Sapthesis come `sapred', ridefinito qui per usarlo già nel preambolo);
%    - cornici dei link di hyperref: rosso [1 0 0] per i riferimenti interni,
%      verde [0 1 0] per le citazioni, ciano [0 1 1] per gli URL.
%  Tutti gli altri colori sono tinte di questi, oppure i colori ufficiali
%  delle classi ISPRS, che servono a leggere maschere e legende.
% =====================================================================

% --- Colori di identità (identici alla tesi del 2022) -----------------
\definecolor{SapRosso}{rgb}{0.5098039,0.1411765,0.2}   % = sapred, #822433
\definecolor{LinkRosso}{rgb}{1,0,0}                    % riferimenti interni
\definecolor{CitaVerde}{rgb}{0,1,0}                    % citazioni
\definecolor{UrlCiano}{rgb}{0,1,1}                     % URL

% --- Tinte di supporto per schemi e tabelle ---------------------------
\colorlet{SapRossoScuro}{SapRosso!80!black}
\colorlet{SapRossoMedio}{SapRosso!60}
\colorlet{SapRossoTenue}{SapRosso!30}
\colorlet{SapRossoChiaro}{SapRosso!12}
\colorlet{SapRossoVelo}{SapRosso!5}
\colorlet{SapGrigio}{black!55}
\colorlet{SapGrigioTenue}{black!25}
\colorlet{SapGrigioChiaro}{black!8}

% --- Classi ISPRS (palette ufficiale delle maschere di Potsdam/Vaihingen)
\definecolor{ClsImpervious}{RGB}{255,255,255}   % superfici impermeabili
\definecolor{ClsBuilding}{RGB}{0,0,255}         % edifici
\definecolor{ClsLowVeg}{RGB}{0,255,255}         % vegetazione bassa
\definecolor{ClsTree}{RGB}{0,255,0}             % alberi
\definecolor{ClsCar}{RGB}{255,255,0}            % auto
\definecolor{ClsClutter}{RGB}{255,0,0}          % clutter / sfondo

% --- Ruoli semantici negli schemi --------------------------------------
\colorlet{ColEncoder}{SapRossoChiaro}     % encoder / backbone
\colorlet{ColDecoder}{SapGrigioChiaro}    % decoder
\colorlet{ColAttn}{SapRossoTenue}         % blocchi di attenzione
\colorlet{ColQuant}{SapRosso}             % quantizzazione vettoriale
\colorlet{ColFreccia}{SapGrigio}          % frecce di flusso
\colorlet{ColGradiente}{SapRosso}         % percorsi del gradiente

% --- Split dei dati -----------------------------------------------------
\colorlet{SplitTrain}{SapRossoMedio}
\colorlet{SplitVal}{SapRossoTenue}
\colorlet{SplitTest}{SapGrigio}
\colorlet{SplitEscl}{white}

% --- Righe evidenziate e segnaposto -------------------------------------
\colorlet{RigaRiferimento}{SapRossoVelo}  % righe di riferimento nelle tabelle
\colorlet{ColTodo}{SapRossoChiaro}        % sfondo delle note \todo
\colorlet{ColTodoBordo}{SapRosso}
